\documentclass{SYSUReport}
\usepackage{tabularx} % 在导言区添加此行
\usepackage{array}
\usepackage{float}
\usepackage{booktabs}
\usepackage{multirow}
\usepackage{graphicx}
\usepackage{caption}
\usepackage{svg}
\usepackage{listings}
\usepackage{xcolor}
\usepackage{fontspec}

\usepackage{tikz}
\usetikzlibrary{arrows.meta, positioning, shapes.geometric, shadows.blur, backgrounds}


% --- 定义颜色 (模仿 VS Code 或常见 IDE) ---
\definecolor{codegreen}{rgb}{0,0.6,0}
\definecolor{codegray}{rgb}{0.5,0.5,0.5}
\definecolor{codepurple}{rgb}{0.58,0,0.82}
\definecolor{backcolour}{rgb}{0.95,0.95,0.92}
\definecolor{codeblue}{rgb}{0.1,0.1,0.9}

% --- 通用设置 ---
\lstset{
    backgroundcolor=\color{backcolour},   
    commentstyle=\color{codegreen},
    keywordstyle=\color{magenta},
    numberstyle=\tiny\color{codegray},
    stringstyle=\color{codepurple},
    basicstyle=\ttfamily\small,          % 使用等宽字体
    breakatwhitespace=false,         
    breaklines=true,                 % 自动折行
    captionpos=b,                    % 标题在下
    keepspaces=true,                 
    numbers=none,                    % 行号在左
    numbersep=5pt,                  
    showspaces=false,                
    showstringspaces=false,
    showtabs=false,                  
    tabsize=4,
    frame=single,                    % 代码框
    rulecolor=\color{black}          % 边框颜色
}

% --- Python 特定样式 ---
\lstdefinestyle{pythonstyle}{
    language=Python,
    morekeywords={self, cls, @property}, % 可以添加额外的关键字
    deletekeywords={list,dict}           % 如果不想高亮某些词可以删除
}

% --- C 特定样式 ---
\lstdefinestyle{cstyle}{
    language=C,
    morekeywords={printf, scanf, uint32_t, uint16_t, size_t}, % 添加 C 常用函数或类型
}
% --- 终端/Shell颜色定义 ---
\definecolor{termbackslash}{rgb}{0.2,0.2,0.2} % 深灰色背景
\definecolor{termwhite}{rgb}{0.94,0.94,0.94}  % 接近白色的文字

% --- Terminal 特定样式 ---
\lstdefinestyle{termstyle}{
    language=bash,
    backgroundcolor=\color{termbackslash}, % 使用深色背景模拟终端
    basicstyle=\ttfamily\small\color{termwhite}, % 白色文字
    keywordstyle=\color{cyan}\bfseries,       % 指令用青色高亮
    commentstyle=\color{gray},
    stringstyle=\color{codepurple},
    numbers=none,                             % 终端通常不需要行号
    showstringspaces=false,
    frame=single,                             % 单线框
    rulecolor=\color{gray},
    breaklines=true,
    % 定义常用的 Linux/开发指令
    morekeywords={gcc, g++, python3, make, cmake, ./, mpirun, mpicc, ls, cd, sudo, pip},
    % 自动高亮提示符 $ 符号
    literate={\$}{{\textcolor{green}{\$}}}1,
}
\headr{多模态大模型原理与应用}
\lessonTitle{多模态大模型原理与应用课程报告}
\reportTitle{基于视觉理解与自动操作的智能浏览器助手}
\stuname{}
\stuid{}
\inst{计算机学院}
\major{计算机科学与技术}
\date{}
\lstset{ 
    language=SQL,
    basicstyle=\ttfamily\small,
    keywordstyle=\color{blue}\bfseries,
    commentstyle=\color{gray},
    stringstyle=\color{red},
    showstringspaces=false,
    breaklines=true,
    numbers=none,
    frame=single
}
\lstdefinelanguage{SQL}{
  morekeywords={SELECT, FROM, WHERE, ORDER, BY, GROUP, HAVING, INSERT, INTO, VALUES,
                UPDATE, SET, DELETE, CREATE, TABLE, DROP, ALTER, AS, JOIN, LEFT,
                RIGHT, INNER, OUTER, ON, LIMIT, OFFSET, DISTINCT, COUNT, SUM,
                AVG, MAX, MIN, LIKE, IN, NOT, NULL, IS, EXISTS, CASE, WHEN, THEN,
                ELSE, END, AND, OR, BETWEEN},
  sensitive=false,
  alsoletter={_},
  morecomment=[l]{--},
  morecomment=[s]{/*}{*/},
  morestring=[d]',
  morestring=[d]"
}
\lstdefinelanguage{Python}{
  morekeywords={%
    and, as, assert, async, await, break, class, continue, def, del, elif,
    else, except, False, finally, for, from, global, if, import, in, is,
    lambda, None, nonlocal, not, or, pass, raise, return, True, try, while,
    with, yield, self, print, len, range, enumerate, zip, map, filter,
    isinstance, type, str, int, float, list, dict, set, tuple, open, close
  },
  sensitive=true,
  morecomment=[l]{\#},
  morestring=[b]",
  morestring=[b]',
  morestring=[s]{"""}{"""},   % 支持三引号字符串
  morestring=[s]{'''}{'''}
}
\lstdefinelanguage{C++}{
  morekeywords={%
    alignas, alignof, and, and_eq, asm, auto, bitand, bitor, bool, break,
    case, catch, char, char8_t, char16_t, char32_t, class, compl, const,
    consteval, constexpr, const_cast, continue, co_await, co_return, co_yield,
    decltype, default, delete, do, double, dynamic_cast, else, enum, explicit,
    export, extern, false, float, for, friend, goto, if, inline, int, long,
    mutable, namespace, new, noexcept, not, not_eq, nullptr, operator, or,
    or_eq, private, protected, public, register, reinterpret_cast, requires,
    return, short, signed, sizeof, static, static_assert, static_cast, struct,
    switch, template, this, thread_local, throw, true, try, typedef, typeid,
    typename, union, unsigned, using, virtual, void, volatile, wchar_t,
    while, xor, xor_eq,
    % 常用标准库函数/类型（可选）
    cout, cin, endl, string, vector, map, set, pair, make_pair, size_t,
    size, push_back, emplace_back, begin, end, find, sort, swap, move
  },
  sensitive=true,
  alsoletter={_},
  morecomment=[l]{//},
  morecomment=[s]{/*}{*/},
  morestring=[b]",
  morestring=[b]',
  morestring=[s]{R"xxx(}{)xxx"} % 支持原始字符串字面量（简化处理，实际需更复杂规则）
}
\lstdefinelanguage{Java}{
  morekeywords={%
    abstract, assert, boolean, break, byte, case, catch, char, class,
    const, continue, default, do, double, else, enum, extends, final,
    finally, float, for, goto, if, implements, import, instanceof, int,
    interface, long, native, new, package, private, protected, public,
    return, short, static, strictfp, super, switch, synchronized, this,
    throw, throws, transient, try, void, volatile, while, false, null, true
  },
  sensitive=true,
  morecomment=[l]{//},
  morecomment=[s]{/*}{*/},
  morestring=[b]",
  morestring=[b]'
}
\newcolumntype{C}[1]{>{\centering\arraybackslash}p{#1}}
\begin{document} %在document环境中撰写文档

%%%%%%%%%%%%%%%%%%%%%%%%%%%%%%%%%%%%%%%%
% 封面及目录，无需改动此处代码
% 面页，需要在导言区用\sysucseset命令填写基本信息
\cover
% 目录
%\tableofcontents
\cleardoublepage
% 设置正文页眉页脚格式，并重新设置页码为1
%%%%%%%%%%%%%%%%%%%%%%%%%%%%%%%%%%%%%%%%
\fontsize{12}{18}\selectfont
\tableofcontents

\subsection{应用背景}
随着互联网的快速发展，网页的样式和信息量都在经历日新月异地增长，为人们完成特定的网页任务带来了挑战。
在多模态大模型技术成熟的背景下，基于VLM/LLM技术实现的多模态网页Agent能够实现辅助和某种程度地替代人类
完成对应的网页任务，具有广阔的应用前景。Web Agent可以作为自动化测试工具，辅助软件开发团队进行前端的
冒烟测试与回归测试。Web Agent也可以作为智能个人助理，帮助个人用户执行复杂任务，降低网页任务的操作成本。
除此以外，Web Agent可以实现复杂数据的收集与分析、改进针对障碍人士的辅助功能等任务。

\subsection{任务定义}
我们计划实现一个能够理解自然语言指令，并通过浏览器自主完成复杂任务的智能代理。
用户将网页任务目标作为输入，在执行任务的每一步中
获取网页的 Accessibility Tree并根据网页截图进行SoM标记，
将网页截图与SoM输入多模态大模型，生成对应的指令，并使用playwright执行并等待页面反馈、进行任务的下一步。
基于任务导向，我们计划实现：
\begin{itemize}
    \item 一个SoM内容标注器，根据网页截图和Accessibility Tree生成元素边界框和对应的id，提高模型的准确率
    \item 一个基于VLM的Web Agent，根据用户输入进行任务的分步计划、在每一步根据网页截图信息和SoM标注生成对应的网页操作指令
    \item 一个playwright执行器，解析Agent输出的执行指令并执行，刷新网页进行下一步任务
\end{itemize}


\section{相关工作调研}
本课题相关工作涉及多模态GUI代理、网页自动化框架、任务规划与视觉定位等方向，以下列举代表性论文与开源项目。

\subsection{多模态网页Agent}
\begin{itemize}
    \item \textbf{WebVoyager} \cite{webvoyager}: 提出基于GPT-4V的端到端网页Agent，在真实网站上通过截图和HTML辅助实现高成功率，验证了大视觉模型在网页导航中的可行性。
    \item \textbf{Fara-7B} (Microsoft Research) \cite{fara7b}: 一个7B参数的多模态模型，结合Magentic-UI框架，使用Playwright执行操作，在控制复杂度的同时实现良好性能，是本课题的核心参考。
    \item \textbf{WebAgent} 综述 \cite{webagent_survey}: 对基于LLM的网页自动化Agent进行了系统分类，总结了规划、感知、执行的典型架构。
\end{itemize}

\subsection{视觉定位与混合感知}
\begin{itemize}
    \item \textbf{SoM (Set-of-Mark)} \cite{som}: 通过为可交互元素叠加数字标记，将 GUI 定位任务转化为标记选择任务，显著提升 VLM 在界面上的 grounding 准确率。本系统将 SoM 作为预处理模块，为 VLLM 提供可直接引用的标记索引。
    \item \textbf{R-VLM} \cite{rvlm}: 提出IoU感知的目标函数和缩放区域候选框，将GUI元素定位精度大幅提升，在Mind2Web基准上获得高准确率，为改进视觉定位提供了理论支持。
    \item \textbf{WebGUM} \cite{webgum}: 证明同时利用HTML/可访问性树和视觉信息比纯视觉或纯DOM方法更有效，为本课题混合感知策略提供了实验依据。
\end{itemize}

\subsection{相关开源项目}
\begin{itemize}
    \item \textbf{Playwright} \cite{playwright}: 现代浏览器自动化框架，提供跨浏览器、可靠的选择器和等待机制，本课题的执行层核心。
    \item \textbf{OpenClaw / Nanobot} \cite{openclaw}: 开源的多模态Agent实现，展示了从截图到操作的简易流水线，但缺乏精心设计的错误恢复和混合感知。
    \item \textbf{SMART-LLM} \cite{smartllm}: 多Agent协作框架，用于复杂任务规划，可借鉴其任务分解与验证机制。
\end{itemize}

\section{系统方案设计}

\subsection{SoM视觉输入模式}

SoM（Set-of-Mark）模式是系统的主要运行模式。该模式将网页截图与可交互元素的视觉标记相结合，使VLM能够直观地看到页面上的可操作元素，并通过引用标记编号来精确指定操作目标。其整体架构如图\ref{fig:arch_som}所示，遵循ReAct（Reasoning + Acting）范式，每步操作分为感知（Perceive）、推理（Reason）、执行（Act）和验证（Verify）四个阶段。

\begin{figure}[htbp]
\centering
\begin{tikzpicture}[
    node distance=0.8cm and 1.5cm,
    block/.style = {
        rectangle, draw, thick, rounded corners=3pt,
        fill=blue!8, drop shadow,
        text centered, font=\sffamily\footnotesize, minimum width=2.4cm, minimum height=0.8cm
    },
    subblock/.style = {
        rectangle, draw, thin, rounded corners=2pt,
        fill=gray!5, font=\sffamily\scriptsize, minimum width=2.2cm, minimum height=0.55cm
    },
    arrow/.style = {thick, -{Stealth}},
    label/.style = {font=\sffamily\scriptsize}
]

% 四阶段主节点
\node (perceive) [block, fill=yellow!10] {1. 感知 (Perceive)};
\node (reason) [block, fill=orange!10, below=0.5cm of perceive] {2. 推理 (Reason)};
\node (act) [block, fill=green!10, below=0.5cm of reason] {3. 执行 (Act)};
\node (verify) [block, fill=red!10, below=0.5cm of act] {4. 验证 (Verify)};

% 感知子模块
\node (extract) [subblock, right=0.3cm of perceive, xshift=2.8cm] {ElementExtractor};
\node (marker) [subblock, below=0.2cm of extract] {SoMMarker};
\node (bridge) [subblock, below=0.2cm of marker] {DOMBridge};

% 推理子模块
\node (vlm_call) [subblock, right=0.3cm of reason, xshift=2.8cm] {VLM调用 (OpenAI/Qwen等)};
\node (parser) [subblock, below=0.2cm of vlm_call] {ActionParser};

% 执行子模块
\node (playwright) [subblock, right=0.3cm of act, xshift=2.8cm] {Playwright浏览器};
\node (executor) [subblock, below=0.2cm of playwright] {ActionExecutor (4层回退)};

% 验证子模块
\node (pixel_diff) [subblock, right=0.3cm of verify, xshift=2.8cm] {像素对比 (全图)};
\node (dom_check) [subblock, below=0.2cm of pixel_diff] {DOM错误检测};
\node (vlm_verify) [subblock, below=0.2cm of dom_check] {VLM验证 (可选)};

% 阶段间主流程
\draw[arrow] (perceive) -- (reason);
\draw[arrow] (reason) -- (act);
\draw[arrow] (act) -- (verify);
\draw[arrow, dashed, blue] (verify.west) -- ++(-1.2,0) |- (perceive.west)
    node[pos=0.5, left, label] {下一轮};

% 输入输出
\node (input_label) [font=\sffamily\tiny, text=gray, left=0cm of perceive, xshift=-2.0cm] {任务描述 + URL};
\node (output_label) [font=\sffamily\tiny, text=gray, left=0cm of verify, xshift=-2.0cm] {Answer / 操作};
\draw[arrow, gray!50] (input_label) -- (perceive);
\draw[arrow, gray!50] (verify) -- (output_label);

\end{tikzpicture}
\caption{SoM模式的ReAct循环架构}
\label{fig:arch_som}
\end{figure}

\textbf{感知阶段（Perceive）}负责将当前网页状态转化为VLM可理解的表示。首先通过Playwright截取当前页面的全页面截图，同时通过\texttt{ElementExtractor}模块使用CSS选择器（覆盖\texttt{<button>}、\texttt{<a>}、\texttt{<input>}等原生交互标签，以及\texttt{[role]}、\texttt{[aria-label]}、\texttt{[data-a11y]}等富交互属性）提取页面中所有可交互的DOM元素。对于现代SPA框架中无任何HTML属性的按钮，通过检测\texttt{cursor: pointer}计算样式作为补充识别手段。元素经过可见性、尺寸（最小4px²）、disabled和pointer-events过滤后，由\texttt{SoMMarker}模块在截图上绘制带编号的彩色矩形边框（12色调色板）。同时\texttt{DOMBridge}维护SoM编号到CSS选择器及Playwright定位器的双向映射，使后续执行阶段能将VLM输出的元素编号解析为可点击的DOM元素。对于旧版QQ邮箱等使用frameset架构的网页，感知阶段会遍历所有\texttt{<iframe>}子文档并修正坐标偏移，确保标注位置正确。

\textbf{推理阶段（Reason）}将SoM标注后的截图以base64 JPEG格式连同任务描述、当前页面URL/标题、历史操作记录（含验证成功/失败标记）和可交互元素文字列表（编号、标签、文字内容、aria-label）一并发送给VLM。VLM返回的结构化JSON包含操作类型（\texttt{click/type/scroll/press/goto/answer/fail}）、目标元素编号和操作值，同时还包含对已完成步骤的总结和后续计划，帮助模型在跨步骤任务中保持连贯性。

\textbf{执行阶段（Act）}由\texttt{ActionExecutor}负责将VLM决策转化为浏览器操作。点击操作采用4层回退机制：优先使用Playwright原生点击（产生\texttt{isTrusted=true}的真实事件），失败后依次尝试JS MouseEvent派发、基于元素真实链接的\texttt{goto}导航（绕过Bing等网站的点击拦截）、坐标中心点的\texttt{elementFromPoint}命中，最后回退至屏幕坐标点击。对于TYPE操作，强制要求先点击输入框获取焦点后再输入。执行完成后检测URL变化和新增标签页，自动切换到最新打开的标签页。

\textbf{验证阶段（Verify）}对操作效果进行判断。首先通过PIL的\texttt{ImageChops.difference}进行全图像素级对比，差异像素少于$0.05\%$时直接判定为无效操作。随后扫描DOM中的错误弹窗和成功提示（QQ邮箱的xmail-ui-dialog等）。若以上检测均未触发，可选调用VLM进行视觉对比验证。验证失败的步骤在下一轮推理时以\texttt{[NO EFFECT]}标记显示在历史记录中，引导模型尝试替代方案。

\subsection{HTML文本输入模式}

HTML模式是为Mind2Web数据集评测而引入的纯文本输入变体。其设计目标是让同一个Agent框架能够在不依赖截图的情况下，仅通过HTML元素文本列表完成网页操作评估。由于Mind2Web数据集预先标注了每个步骤的正负候选元素及其\texttt{backend\_node\_id}，HTML模式不需要截图和SoM标注，改为直接从数据集读取候选列表并格式化为编号文本供LLM选择。

HTML模式与SoM模式共享完全相同的Agent循环和Executor，差异仅在于Perceptor和Prompt。如图\ref{fig:html_mode}所示：

\begin{itemize}
    \item \textbf{感知层}：由\texttt{HTMLPerceptor}替代\texttt{SoMPerceptor}，离线评估时读取Mind2Web候选数据，实时运行时可回退至\texttt{ElementExtractor}提取真实页面元素，同时提取\texttt{body.innerText}作为页面内容摘要供LLM理解当前页面状态。
    \item \textbf{推理层}：Reasoner自动检测\texttt{perception.screenshot}是否为空，若为空则切换至HTML文本Prompt，将元素列表、页面文本、历史记录和当前计划格式化为纯文本发送给LLM。
    \item \textbf{验证层}：HTML模式下无截图，验证功能被简化为URL变化检测和DOM错误弹窗扫描。
\end{itemize}

\begin{figure}[htbp]
\centering
\begin{tikzpicture}[
    node distance=0.6cm and 2cm,
    block/.style = {
        rectangle, draw, thick, rounded corners=3pt,
        fill=blue!6, drop shadow,
        text centered, font=\sffamily\footnotesize, minimum width=2.6cm, minimum height=0.8cm
    },
    somblock/.style = {block, fill=yellow!6, align=center},
    htmlblock/.style = {block, fill=green!6, align=center},
    arrow/.style = {thick, -{Stealth}},
    dasharrow/.style = {thick, -{Stealth}, dashed, gray},
]

% SoM侧
\node (som_label) [font=\sffamily\footnotesize\bfseries, text=blue] at (0, 0.2) {SoM视觉模式};
\node (som_p) [somblock, below=0.3cm of som_label] {SoMPerceptor\\截图+标注+提取};
\node (som_r) [somblock, below=0.3cm of som_p] {SoM Prompt\\截图+元素列表+历史};

% HTML侧
\node (html_label) [font=\sffamily\footnotesize\bfseries, text=green!50!black, right=4.5cm of som_label] {HTML文本模式};
\node (html_p) [htmlblock, below=0.3cm of html_label] {HTMLPerceptor\\候选/提取+页面文本};
\node (html_r) [htmlblock, below=0.3cm of html_p] {HTML Prompt\\元素文本+页面文本+历史};

% 共享层
\node (shared) [block, fill=gray!10, minimum width=7.5cm] at (2.25,-3.2) {共享: Agent循环 / ActionExecutor / Playwright};

% 箭头
\draw[arrow] (som_p) -- (som_r);
\draw[arrow] (html_p) -- (html_r);
\draw[dasharrow] (som_r.south) -- ++(0,-0.3) -- ++(-1.5,0) -- (shared.north -| som_r);
\draw[dasharrow] (html_r.south) -- ++(0,-0.3) -- ++(1.5,0) -- (shared.north -| html_r);
\draw[dasharrow, gray] (som_p.east) -- ++(0.5,0) |- (html_p.west)
    node[pos=0.5, above, font=\sffamily\tiny, text=gray] {ElementExtractor可复用};

\end{tikzpicture}
\caption{SoM模式与HTML模式的架构对比}
\label{fig:html_mode}
\end{figure}

HTML模式无需额外开发，仅需替换Perceptor和Prompt即可从视觉输入切换为文本输入，有效降低了为不同数据集定制评测流程的成本。

\section{实验评测与分析}
针对项目模型的SoM输入与HTML输入模式分别设计了评测实验。所有本地模型均在一台RTX 3090 24GB $\times 1$，内存60GB的计算机上运行，并使用Qwen 3VL-8B-Instruct作为本地测试模型,使用API调用Qwen3.7-max测试。
选取任务成功率(SR)、元素准确率(Ele. Acc)、任务单步操作的F1分数(OP.F1)、任务单步成功率(Step SR)以及单步平均耗时作为评价Web Agent性能的标准。
其中，Ele. Acc表示各任务内选择正确网页元素的步数比例跨任务的均值，Op.F1分数表示单步操作的F1值，Step SR表示各任务内元素与操作均正确的步数比例跨任务的均值，SR表示所有步骤均成功执行的任务数占数据集总任务数的比例。

\begin{align*}
    \text{Ele.Acc} &= \frac{1}{M} \sum_{i=1}^{M} \frac{1}{N_i} \sum_{j=1}^{N_i} \mathbb{I}(\hat{e}_{i,j} \in S_{acc}(e_{i,j})) \\[6pt]
    \text{Op.F1}_{i,j} &= \mathbb{I}\bigl(\hat{op}_{i,j} \in S_{op}(op_{i,j})\bigr) \times
                         \begin{cases}
                            1 & op_{i,j} = \text{CLICK} \\[4pt]
                            \text{F1}(v_{i,j}, \hat{v}_{i,j}) & op_{i,j} \in \{\text{TYPE}, \text{SELECT}\}
                         \end{cases} \\[4pt]
    \text{Step SR} &= \frac{1}{M} \sum_{i=1}^{M} \frac{1}{N_i} \sum_{j=1}^{N_i} \Bigl[ \mathbb{I}(\hat{e}_{i,j} \in S_{acc}(e_{i,j})) \times \mathbb{I}(\text{Op.F1}_{i,j} = 1) \Bigr] \\[6pt]
    \text{SR}_i &= \prod_{j=1}^{N_i} \Bigl[ \mathbb{I}(\hat{e}_{i,j} \in S_{acc}(e_{i,j})) \times \mathbb{I}(\text{Op.F1}_{i,j} = 1) \Bigr]
\end{align*}
其中$M$为数据集内的总任务数，$N_{i}$为第$i$个任务的步骤数，$e_{i,j},\hat{e}_{i,j}$分别为第$i$个任务第$j$步中Ground Truth的等价网页元素集合和Agent预测的网页元素，$op_{i,j} \in \{\text{CLICK},\text{TYPE},\text{SELECT}\}$为Ground Truth操作类型，$\hat{op}_{i,j}$为预测操作类型，$v_{i,j},\hat{v}_{i,j}$分别为Ground Truth和预测的操作值。
$S_{op}(\text{CLICK})=\{\text{CLICK}, \text{HOVER}, \text{PRESS}\}$为Mind2Web中将$\text{HOVER}$和$\text{PRESS\_ENTER}$映射至$\text{CLICK}$的等价操作集合；$\text{F1}(\cdot,\cdot)$为F1分数。

\subsection{HTML输入模式下的性能测评}
选用Mind2Web数据集对模型进行评测。Mind2Web数据集是一个大规模的Web Agent基准数据集,将真实世界的静态网页作为数据,具有1,341个不同任务和46步的最长任务步数，在Web Agent领域被广泛用于测试模型的页面理解与执行能力。

我们对Mind2Web数据集提供的评测方法进行了修改。在Mind2Web数据集提供的HTML应用过滤规则只保留可交互元素后，使用一个基于Bert2实现的预训练模型all-MiniLM-L6-v2对剩余候选元素进行排序并交给VLM进行选择，并给出对应的操作类型和元素，最后计算Ele. Acc,OP.F1,Step SR和SR。
\begin{table}[H]
\centering
\small
\caption{HTML输入模式下模型在Mind2Web数据集上的性能}
\label{tab:html_perf}
\setlength{\tabcolsep}{3pt}
\begin{tabular}{l|rrrr|rrrr|rrrr|r}
\toprule
& \multicolumn{4}{c}{Cross-Task} & \multicolumn{4}{c}{Cross-Website} & \multicolumn{4}{c}{Cross-Domain} & \\
\cmidrule(lr){2-5} \cmidrule(lr){6-9} \cmidrule(lr){10-13}
& EA & OP & SS & SR & EA & OP & SS & SR & EA & OP & SS & SR & Time(s) \\
\midrule
Qwen3-VL-8B & 41.9 & 84.0 & 33.9 & 1.5 & 32.0 & 78.0 & 26.4 & 1.0 & 36.0 & 83.0 & 28.0 & 3.0 & 1.3 \\
Qwen3.7-max & 64.6 & 80.5 & 52.0 & 8.7 & 57.9 & 81.3 & 42.4 & 6.2 & 55.6 & 81.9 & 43.8 & 7.3 & 14.3 \\
\bottomrule
\end{tabular}
\end{table}
由表1可以看出，Qwen3.7-max在元素定位上显著优于Qwen3-VL-8B，元素准确率平均比Qwen3-VL-8B提高了$22.7\%$，最终导致更高的单步成功率与任务成功率。但Qwen3-VL-8B作为小模型具有快速响应与低成本
的优势，在相对较好的网页操作性能和较快的响应时间中取得了平衡。Qwen3.7-max的参数量过大，虽然准确率更高但响应时间是Qwen3-VL-8B的11倍，不具有实时性。
\subsection{SoM输入模式下的性能测评}

为评估SoM视觉输入模式在真实网页场景下的表现，自行设计了包含30个任务的测试集，全部基于可公开访问的真实网站。测试任务覆盖多种难度与应用场景，代表性任务包括：搜索"中山大学广州校区南校园"并返回地址、打开哔哩哔哩首页点击第一个视频后返回作者头像、登录QQ邮箱并向指定地址发送端午祝福邮件等。完整任务列表见附录。所有测试均使用Qwen 3VL-8B-Instruct模型，在有头浏览器模式下运行，每个任务至少重复运行10次以减小随机性。

由于自建测试集缺少Mind2Web数据集中的Ground Truth标注（如真实操作序列和正负候选元素），无法直接计算Ele. Acc和Op.F1等微观指标。因此针对真实网页场景，定义以下两项宏观指标评价Agent的任务完成能力：

\begin{align*}
    \text{SR} &= \frac{S}{N} \\
    \text{AvgSteps} &= \frac{1}{S} \sum_{i \in \mathcal{S}} \text{steps}_i
\end{align*}
其中$N=30$为测试任务总数，$S$为成功完成（Agent返回正确答案且人工验证通过）的任务数，$\mathcal{S}$为成功任务集合，$\text{steps}_i$为第$i$个任务成功完成时使用的步数。

\begin{table}[H]
\centering
\small
\caption{SoM输入模式与HTML输入模式在自建测试集上的对比}
\label{tab:som_vs_html}
\setlength{\tabcolsep}{6pt}
\begin{tabular}{l|rr|rr}
\toprule
\multirow{2}{*}{模型} & \multicolumn{2}{c|}{SoM输入} & \multicolumn{2}{c}{HTML输入} \\
\cmidrule(lr){2-3} \cmidrule(lr){4-5}
& SR & AvgSteps & SR & AvgSteps \\
\midrule
Qwen 3VL-8B & 73.3 & 10.5 & 40.0 & 6.8 \\
\bottomrule
\end{tabular}
\end{table}

由表2可以看出，SoM视觉输入模式相比HTML文本输入模式在真实网页任务上具有显著优势。任务成功率从$40.0\%$提升至$73.3\%$，提升幅度达$83.3\%$。分析任务完成情况发现，HTML输入模式能够成功完成的任务以简单的搜索引擎查询为主（如事实类搜索、天气查询等），一旦任务涉及表单填写、页面跳转、多步交互等复杂操作，成功率显著下降。SoM模式由于能够直观地识别页面布局与元素空间位置，在需要点击特定按钮、填写表单、切换标签页等复杂操作时表现更好。但SoM模式的完成任务平均步数$10.5$步明显高于HTML模式的$6.8$步，这一方面是因为SoM模式成功完成的任务本身更复杂，自然需要更多操作步数；另一方面也与SoM模式下模型倾向于执行更多探索性操作（如下拉滚动、尝试不同元素）有关。

总体而言，SoM输入模式以更多的操作步数为代价，换取了更高的任务完成率与更广泛的任务适用性，在真实网页场景下优于HTML文本输入模式。

\section{参考文献}
\begin{thebibliography}{99}
\bibitem{webvoyager} He, Y., et al. “WebVoyager: Building an End-to-End Web Agent with Large Multimodal Models.” 10.48550/arXiv.2401.13919
\bibitem{fara7b} Microsoft Research. “Fara-7B: A Foundation Model for GUI Automation.” 2025.
\bibitem{webagent_survey} Xinyi Li, et al. “A survey on LLM-based multi-agent systems: workflow, infrastructure, and challenges”  Computer Science.10.1007/s44336-024-00009-2
\bibitem{som} Yang, J., et al. “Set-of-Mark Prompting Unleashes Extraordinary Visual Grounding in GPT-4V.” arXiv preprint arXiv:2310.11441, 2023.
\bibitem{rvlm} Cheng, K., et al. “R-VLM: Region-Aware Vision Language Model for Precise GUI Grounding.” arXiv preprint 10.48550/arXiv.2507.05673.
\bibitem{webgum} Furuta, H., et al. “WebGUM: Multimodal Web Navigation with Instruction-Finetuned Foundation Models.” arXiv.2305.11854v4
\bibitem{playwright} Microsoft. “Playwright: Fast and Reliable Cross-Browser Automation.” \url{https://playwright.dev/}, 2025.
\bibitem{openclaw} OpenClaw Contributors. “OpenClaw: Open-Source Multimodal Web Agent.” GitHub, 2025.
\bibitem{smartllm} SMART-Lab Purdue. “SMART-LLM: Multi-Agent Task Planning.” GitHub, 2024.
\end{thebibliography}

\end{document}
